\subsection{El término constante del predictor}\label{sec:D05}
Recordamos la Definición~\ref{def:Qpred}:
\[
 \log Q_{\rm pred}(n)=H(\pi^{(n)})+L_n
           -\tfrac12\log\operatorname{pdet}S_n
           +\tfrac12\log(4\pi n),                    \tag{D5.1}
\]
donde $L_n\to3I_2$ por (LC4).
Aquí $\pi=3{,}14159\ldots$ sólo en el logaritmo de normalización y en la fórmula de Stirling.
Como en la Sección~\ref{sec:preliminares}, $\Phi=s_1+1$, donde $s_1$ es el valor óptimo en
masa uno.
El teorema siguiente es el Teorema~\ref{thm:intro-constante} de la introducción: $J$ es el operador
de~(D3.6), $\mathcal D(J)=\log\det_2(I+J)$ por~(D4.1), e $I+J\ge aI$ con $a>0$ por la
Sección~\ref{sec:D04}.

\begin{teorema}[Término constante del predictor]\label{thm:constante}
Cuando $n\to\infty$,
\[
 \log Q_{\rm pred}(n)-\log n!-n\Phi\longrightarrow C_{\rm pred},
                                                        \tag{D5.2}
\]
\[
 C_{\rm pred}=\tfrac12\log2-\tfrac12
             -\tfrac32 I_2-\tfrac12\mathcal D(J).     \tag{D5.3}
\]
\end{teorema}
\begin{proof}
Reunimos las tres expansiones necesarias, manteniendo en cada línea los términos de orden uno:
\[
\begin{aligned}
 H(\pi^{(n)}) &= n\log n+ns_1+\tfrac32 I_2+o(1) &&\text{por (D2.5)},\\
 L_n &= 3I_2+o(1) &&\text{por (LC4)},\\
 \log\operatorname{pdet}S_n &=1+12I_2+\mathcal D(J)+o(1) &&\text{por (D4.3)}.
\end{aligned}
\]
Por Stirling,
\[
 \log n!=n\log n-n+\tfrac12\log(2\pi n)+o(1).
\]
Al sustituir en (D5.1), los términos principales se cancelan porque
$\Phi=s_1+1$. La normalización deja $\tfrac12\log2$,
y $\tfrac32I_2+3I_2-6I_2=-\tfrac32I_2$.
Así se obtiene (D5.3).
\end{proof}
